% =============================================================================
% Response to Reviewers -- BMC Bioinformatics
% Manuscript: GaugeFixer
%
% Formatting convention (matches the original Word document):
%   black text          -- editor and reviewer critiques (verbatim)
%   responseblue text   -- the authors' responses
%   msquote environment -- excerpts from the revised manuscript (indented, blue)
% =============================================================================
\documentclass[11pt]{article}

\usepackage[margin=1in]{geometry}
\usepackage[T1]{fontenc}
\usepackage[utf8]{inputenc}
\usepackage{lmodern}
\usepackage{amsmath,amssymb}
\usepackage{xcolor}
\usepackage{enumitem}
\usepackage{parskip}
\usepackage[round]{natbib}
\usepackage[hidelinks]{hyperref}

\definecolor{responseblue}{HTML}{0070C0}
\definecolor{linkblue}{HTML}{1155CC}

% Authors' response: blue body text.
\newenvironment{response}{\par\color{responseblue}}{\par}

% Excerpt from the revised manuscript: indented block, same blue.
\newenvironment{msquote}
  {\begin{quote}\color{responseblue}\itshape\noindent\ignorespaces}
  {\end{quote}}

% Reviewer/editor comments are numbered items with generous separation.
\setlist[enumerate,1]{leftmargin=*,itemsep=1.2\baselineskip,topsep=\baselineskip}

\newcommand{\reviewersection}[1]{%
  \bigskip\par\noindent{\large\bfseries #1}\par\medskip\hrule\medskip}

\begin{document}

% =============================================================================
\begin{center}
{\Large\bfseries Response to Reviewers}\\[0.6em]
{\large GaugeFixer: overcoming parameter non-identifiability in models of\\
sequence-function relationships}\\[0.8em]
Carlos Mart\'{i}-G\'{o}mez, David M.\ McCandlish, Justin B.\ Kinney\\[0.5em]
Submission ID: 9a7ff349-d9aa-48f2-8210-3e1125bbde71\\
7 August 2026
\end{center}

\bigskip

\begin{response}
We thank the Editor and the two Reviewers for their consideration of our manuscript. This work has been substantially improved as a result of the feedback we recieved. Additional edits have been made throughout the manuscript to improve clarity. Below is a point-by-point response to each of the critiques raised. 

We also note that, since submitting the initial manuscript, we realized that several methods
proposed by other groups for analyzing sequence-function relationships can
also be computed efficiently using the algorithms we describe
\citep{Weinberger1991,Poelwijk2016,Stormo:2011jb,Brookes2022,Tsui2026,Faure:2024gy}.
We have implemented these methods in GaugeFixer, and have added a
paragraph to the Discussion explaining how our algorithms accelerate these
computations and extend them to hierarchical models.
\end{response}

\begin{msquote}
Finally, Algorithms 1 and 2 can also be used for purposes other than gauge
fixing. In particular, the Walsh-Hadamard transformation of sequence-function
landscapes \citep{Weinberger1991, Poelwijk2016} and its generalizations
\citep{Stormo:2011jb, Brookes2022, Tsui2026, Faure:2024gy} involve multiplying
an $\alpha^L$-dimensional vector (representing a combinatorially complete
landscape) by a dense matrix constructed as a Kronecker product of
position-specific matrices. Algorithm 1 allows these transformations to be
performed far more efficiently than by standard matrix multiplication.
Algorithm 2, moreover, extends these transformations to landscapes that are
defined by hierarchical models (e.g., $K$-order or $K$-adjacent models), rather
than by combinatorially complete sets of measurements. This latter capability
should allow these methods to be applied to sequences that are much longer than
is computationally feasible using published approaches. We have incorporated
these other landscape-analysis approaches into GaugeFixer alongside the
package's core gauge-fixing functionality.
\end{msquote}

% =============================================================================
\reviewersection{Editor Comments}

Thank you for submitting your manuscript to BMC Bioinformatics. Please see below
for the comments from the reviewers. The reviewers have raised points that
require minor revision of the manuscript before it is approved for the
publication. Please address their comments carefully.

Dongheon Lee

Please address the points mentioned below along with the editorial comments.

\begin{enumerate}
\item Please ensure that the statement given in the ``Data Availability'' section
of the SNAPP submission system matches the statement included in the
``Availability of Data and Materials'' section of the manuscript. When submitting
your revised manuscript, please select `yes' under the data availability
declaration and input your data availability statement in the appearing textbox.
\end{enumerate}

We recommend submitting all revisions within the mentioned deadline.

If you need more time, please contact us and include your submission ID.

Kind regards,

Payel Karmakar\\
Assistant Editor\\
BMC Bioinformatics

Support contact: \href{mailto:bmcbioinformatics@biomedcentral.com}{bmcbioinformatics@biomedcentral.com}

\medskip

\begin{response}
We have selected ``yes'' under the data availability declaration, and
have entered the manuscript's ``Availability of data and materials''
section into the SNAPP submission system.


\end{response}

% =============================================================================
\reviewersection{Reviewer 1}

In this paper, building on previous work of a more theoretical nature, the authors
define and illustrate a practical method for fixing the ``gauge symmetries'' that
hamper the interpretability of even the most basic sequence-to-phenotype models.
This is an important topic that is not on most people's radar screens: two
equivalent models can have very different parameters even if their predictions
are identical. The authors include a software package that allows users to
standardize the parameters of their own models. This is a very useful
contribution to the field in my opinion.

I have one suggestion: the fact that the boundary case at position $-$9 is
distinct from more upstream positions could have two explanations: (i) an
intrinsically positional effect or (ii) the pentamer window used is not wide
enough to fully capture the sequence-to-phenotype relationship, which, if the
fixed base at position $-$4 in the library design is suboptimal, would partly
explain the steep drop in panel 3B. This would be easy to check explicitly, and
would add to the analysis.

\begin{response}
We thank the reviewer for this suggestion. To distinguish these two
explanations, we repeated the analysis using hierarchical gauges defined by the
extended motifs GAGGAG and AGGAGU, which place the fixed $-14$G and $-4$U of the
experimental construct immediately adjacent to the core motif. The results
support the reviewer's hypothesis in part: the fixed $-4$U does appear to
contribute to the low mean activity of sequences carrying the core motif at
register $-9$, although this contribution is itself register-dependent. The fixed
$-14$G, by contrast, does not account for the reduced activity at register $-13$.
We have updated Figure 3B to include these results and have added the following
text to the Results section:
\end{response}

\begin{msquote}
The dependence of the constant parameter on register could reflect the intrinsic
positional preference of the core motif, the influence of the fixed sequence
flanking the variable region, or both. Every assayed sequence carries a fixed G
at position $-14$ and a fixed U at position $-4$, just outside the variable
region. The effects of these characters cannot be read directly from the model.
They can, however, be probed indirectly: extending the core motif to GAGGAG or
AGGAGU places a G or a U immediately adjacent to the core, and the resulting
constant parameter reports the mean activity of sequences carrying that character
at that position. We therefore computed constant parameters under hierarchical
gauges defined by these extended motifs at every register where they fit within
the variable region (Figure 3B). Neither extended motif can be placed at the
register where its added character would coincide with the corresponding fixed
position in the assay: AGGAGU does not fit at register $-9$, nor GAGGAG at
register $-13$. The neighboring registers, $-10$ and $-12$ respectively,
therefore provide the closest available proxies. At register $-10$, the mean
activity of sequences containing AGGAGU was substantially lower than that of
sequences containing AGGAG alone; a smaller decrease was observed at register
$-13$, and essentially none at registers $-12$ and $-11$. This suggests that the
low average activity of sequences carrying the core motif at register $-9$ is
partly attributable to the negative contribution of the fixed $-4$U, although
this contribution is itself register-dependent. By contrast, the extended motif
GAGGAG raised mean activity at registers $-11$ through $-9$, but not at register
$-12$. The fixed $-14$G is therefore unlikely to explain why the constant
parameter is lower at register $-13$ than at registers $-12$ and $-11$; an
additional upstream G appears to be beneficial only when the core motif lies
closer to the start codon.
\end{msquote}

\medskip
\noindent Minor comments:

\begin{enumerate}

% ---------------------------------------------------------------- R1.1
\item Abstract, ``scales quadratically'': This is in terms of ``$M$'', which
itself is defined as $(1+\alpha)^L$, so exponential in the alphabet size. It
would be good to clarify this.

\begin{response}
The revised abstract clarifies that the size of the projection matrix scales quadratically
with the number of model parameters. 
\end{response}

\begin{msquote}
The most direct computational implementation of these methods is often impractical, however, because it requires a projection matrix whose size scales quadratically with the number of model parameters.
\end{msquote}

\begin{response}
The number of parameters is the relevant
quantity here, because its scaling with sequence length and alphabet size depends
on the specific model considered. Although the total number of parameters is
$(\alpha+1)^L$ for an all-order model, this expression does not apply to the more
general class of hierarchical models. For example, the number of parameters
scales linearly with $L$ in an additive model and quadratically with $L$ in a
pairwise-interaction model. We have revised the Technical Background section to
clarify this distinction:
\end{response}

\begin{msquote}
It is therefore common to use models that include only interactions up to a
specified order ($K$-order models) or only interactions between nearby positions
($K$-adjacent models). The number of parameters in these models grows only
polynomially with sequence length (as $L^K$ for $K$-order models and linearly in
$L$ for $K$-adjacent models), making them tractable for much longer sequences.
\end{msquote}

% ---------------------------------------------------------------- R1.2
\item Figure 1: it is a little unclear to me how the lines runs relative to the
blue plane. Perhaps choosing a different 3D viewpoint would help, and perhaps
other improvements can be made to make this important figure as clear as
possible.

\begin{response}
We have reworked Figure 1 to more
clearly show the positions of the parameter choices, the zero-sum gauge plane, and
the line of equivalent parameters in 3D space.
\end{response}

% ---------------------------------------------------------------- R1.3
\item Page 3, left column, 2nd paragraph, ``restricted to subsets of positions'':
Can the authors clarify whether these subsets can be overlapping or not?

\begin{response}
The subsets of positions used in the decomposition (the generating sets) may
overlap, and generally do so except in special cases such as additive models or
models defined by disjoint generating sets. Such overlaps affect neither the
applicability nor the correctness of the algorithm. We have clarified this point
in the text by adding:
\end{response}

\begin{msquote}
Note that there are often many valid ways to carry out this decomposition due to the generating sets in $W$ overlapping. Because $P$ is a linear operator, however, the result is independent of which decomposition is used.
\end{msquote}

% ---------------------------------------------------------------- R1.4
\item Near Equation (4) the authors also use the terms ``feature'' and ``subset of
positions'' interchangeably. Good to explicitly say that they are the same thing,

\begin{response}
These two terms are actually intended to represent different objects: a feature $x^u_U$ is defined by a subset of positions
$U$ and a subsequence $u$ of length $|U|$. It takes value 1 if the subsequence at
positions $U$ for sequence $s$ matches $u$ exactly, and 0 otherwise. We have
rewritten this passage to make the distinction explicit:
\end{response}

\begin{msquote}
Formally, each feature $x^u_U(s)$ is specified by a subset of positions
$U \subseteq S$ and a subsequence $u \in \mathcal{A}^{|U|}$:
\begin{equation*}
x^u_U(s) = \left\{
  \begin{array}{cl}
    1 & \textrm{if~} s[U] = u \\
    0 & \textrm{otherwise}
  \end{array} \right.,
\end{equation*}
where $s[U]$ indicates the subsequence of $s$ that occurs at the positions $U$.
\end{msquote}

% ---------------------------------------------------------------- R1.5
\item Under Equation 9: Make explicit that since $M = (1+\alpha)^L$, scaling as
$O(L(1+\alpha)M)$ is the same as $O(L(1+\alpha)^{1+L})$, etc.

\begin{response}
We have substantially clarified the scaling relations in the revised manuscript. For all-order models,
$M=(1+\alpha)^L$, so $L=\log_{1+\alpha} M$. Therefore, for a fixed alphabet size
$\alpha$, the computational complexity $O(L(1+\alpha)M)$ is actually equivalent to
$O(M \log M)$, not $O(M)$ as we previously stated. We have revised the text to make this equivalence explicit and to
state the corresponding comparison with direct matrix multiplication:
\end{response}

\begin{msquote}
For each choice of $l$ and indices $i_1, \ldots, i_L$, it takes $\alpha + 1$
multiplication-addition operations to compute the right-hand side of Eq.~10. As
there are $L$ values of $l$ and $(\alpha+1)^L$ possible combinations of indices,
the algorithm has computational complexity $O(L(\alpha+1)^{L+1})$. In terms of
$M$ and $\alpha$, this is $O(\frac{\alpha+1}{\log (\alpha+1)} M \log M)$, which
for fixed $\alpha$ becomes $O(M \log M)$. By contrast, standard matrix multiplication by $P$ requires $O(M^2)$ operations.
\end{msquote}

\begin{response}
The memory requirements are treated in the following paragraph:
\end{response}

\begin{msquote}
Since $M \geq (\alpha+1)^2$ whenever $L \geq 2$, and the algorithm never stores
more than two $M$-sized vectors at once, these requirements scale as $O(M)$. By
contrast, storing $P$ itself requires storing $O(M^2)$ numbers in memory.
\end{msquote}

\begin{response}
We have also added the computational and memory complexities of Algorithm 2:
\end{response}

\begin{msquote}
The memory requirements of this algorithm scale as $O(M)$, just as for
Algorithm 1. The computational requirements, however, depend on the size of $W$
and on the size of each $S' \in W$, both of which are model-specific. For
$K$-order models specifically, the generating sets consist of all possible
subsets of $K$ positions; this means that $|W| = {L \choose K}$ and $|S'| = K$
for all $S' \in W$. The primary computational expense occurs in line 5, which is
carried out ${L \choose K}$ times. Each execution requires $K (\alpha+1)^{K+1}$
operations, following the analysis above. By comparison, the number of parameters is $M = \sum_{k=0}^K {L \choose k}\alpha^k$. Assuming $K \ll L$, we can approximate $M \approx {L \choose K} \alpha^K$. Holding $\alpha$ and $K$ fixed while increasing $L$ thus causes the computational complexity to scale as $O(M)$. $K$-adjacent models behave similarly.
\end{msquote}

% ---------------------------------------------------------------- R1.6
\item Figure 3C: Consider moving label ``register $-$13'' etc from left of heat
map to above the heat map, or even in between rows C and D because it applies to
both.

\begin{response}
The register labels were intended to apply jointly to the two aligned heatmaps,
but we agree that this was not obvious. To make this clearer, we have
relabeled the former panels C and D as the left and right parts of a single panel
C, representing the additive and pairwise parameters for each register,
respectively. We have also adjusted the register labels to clarify that each
applies to both the rectangular and triangular heatmaps in that row.
\end{response}

\end{enumerate}

% =============================================================================
\reviewersection{Reviewer 2}

I find this work of high quality and broad interest and I have only minor
remarks.

\begin{enumerate}

% ---------------------------------------------------------------- R2.1
\item Introduction: Sentence `This non-identifiability occurs because the binary
feature vectors representing the sequences do not span the full space in which
they live' -- could the authors check again this statement, or unpack it? I am not
sure it's true, in the sense that even if I had an exhaustive dataset the issue
has to do with statistical observables (i.e.\ frequencies) not providing the same
amount of constraints as the parameters to infer, as explained at page 5 here:
\href{https://arxiv.org/pdf/1703.01222}{https://arxiv.org/pdf/1703.01222}

\begin{response}
Having fewer independent observables than raw parameters is another way of
describing the same underlying redundancy: generalized one-hot models are
overparametrized, so multiple parameter vectors can represent exactly the same
sequence-function map. Importantly, this is a structural property that results
from the mathematical form of the model, rather than a consequence of the data
used for inference. It therefore applies to any sequence-function model expressed
in this form, including cases in which the function value is known exactly for
every possible sequence. We have revised the text to make this point and the corresponding geometry more explicit.
\end{response}

\begin{msquote}
A fundamental challenge in interpreting the parameters of these models is that,
except in trivial cases, parameter values are not uniquely determined by the
sequence-function landscape that the model describes~\citep{Posfai2025bb}. This
structural non-identifiability is a consequence of the model's mathematical form
and is independent of the type and amount of data used to fit it. It occurs
because the binary feature vectors representing all possible sequences span only
a proper subspace of the vector space in which they are defined. Consequently, a
continuum of parameter vectors yields identical model predictions. These
parameter vectors differ from one another along what are called ``gauge
freedoms''---directions in parameter space that are orthogonal to the binary
feature vectors of all possible sequences.
\end{msquote}

% ---------------------------------------------------------------- R2.2
\item page 3: `In this case there are $M = (\alpha + 1)^L$ parameters' -- just for
clarity, could the author explain why $\alpha+1$ per position, looping back to the
minimal example? (I guess the `+1' is the offset $\theta_0$ there).

\begin{response}
Summing over all subsets of positions directly gives $M=\sum_{k=0}^{L} \binom{L}{k}\alpha^{k}=(\alpha+1)^L$. Alternatively, we show that each
parameter can be indexed by a sequence of length $L$ constructed from an extended
alphabet of $\alpha+1$ characters, the extra character being a wild-card `*'
marking positions not included in $U$. This latter rationale is
actually the basis for the indexing scheme used by the projection matrices.
\end{response}

\begin{msquote}
An all-order model has $M = (\alpha + 1)^L$ parameters. To see this, note that
there are $\binom{L}{k}$ ways to choose a set of positions $U$ of size $k$, and
$\alpha^k$ possible subsequences $u$ for each such set. Summing over $k$ gives
$\sum_{k=0}^{L} \binom{L}{k}\alpha^{k} = (\alpha + 1)^L$ parameters.
Alternatively, we can derive this parameter count by assigning to each feature
$x^u_U$ an index sequence of length $L$ that has subsequence $u$ at the positions
in $U$ and a wild-card character $*$ at all other positions. Because these index
sequences are constructed from the alphabet $\{*, c_1, \ldots, c_\alpha\}$, which
has size $\alpha + 1$, there are $(\alpha+1)^L$ possible index sequences, and hence the
same number of parameters. This indexing approach underlies the tensor
representation introduced by \citet{Posfai2025aa} and used below.
\end{msquote}

% ---------------------------------------------------------------- R2.3
\item Eqs 8-9: I would clarify the range of variation of $i_1, \ldots, i_L$, and
$j_1, \ldots, j_L$

\begin{response}
We have rewritten these equations (now Eqs.~9 and 10) to show explicitly the
range of summation for each index, as well as the fact that all possible
combinations of indices are summed over:
\end{response}

\begin{msquote}
\begin{equation*}
[\vec{\theta}_{\textrm{fixed}}]_{i_1 \cdots i_L} =
\sum_{j_1=1}^{\alpha+1}\sum_{j_2=1}^{\alpha+1} \cdots
\sum_{j_L=1}^{\alpha+1} [P_1]_{i_1 j_1} \cdots [P_L]_{i_L j_L}
[\vec{\theta}]_{j_1 \cdots j_L},
\end{equation*}
where each index $i_l$ runs from $1$ to $\alpha+1$.
\end{msquote}

% ---------------------------------------------------------------- R2.4
\item `To project a hierarchical model into a hierarchical gauge' $\rightarrow$
`onto'?

\begin{response}
We have rewritten this sentence, which now reads:
\end{response}

\begin{msquote}
To express the parameters of a hierarchical model in a hierarchical gauge, we use
the fact that these parameters can be decomposed into a sum of the parameters of
multiple all-order models, each restricted to one of the generating position
sets.
\end{msquote}

% ---------------------------------------------------------------- R2.5
\item Page 4: around algorithm 2, I was a bit confused as to why in some instances
one has $\vec{\theta}[S']$ and in others $\vec{\theta}_{S'}[S']$, do they mean the
same thing? Can one homogenize or clarify the notation?

\begin{response}
We apologize for the confusion. $\vec{\theta}_{S'}[S']$ and $\vec{\theta}_{S'}$
convey the same information but differ in their dimensions: $\vec{\theta}_{S'}$ has
$(\alpha+1)^L$ entries, of which only $(\alpha+1)^{|S'|}$ can be nonzero;
$\vec{\theta}_{S'}[S']$ is the $(\alpha+1)^{|S'|}$-dimensional vector restricted to
those nonzero entries. We agree that this notation is cryptic, and have
therefore clarified Algorithm 2 by introducing an intermediate variable
$\vec{\omega} \gets \vec{\theta}_{\textrm{tmp}}[S']$. We have also clarified the
meaning of the bracket notation:
\end{response}

\begin{msquote}
Here, $\vec{\theta}[S']$ indicates an $(\alpha+1)^{|S'|}$-dimensional vector
composed only of the elements $\theta_U^u$ for which $U \subseteq S'$, and
$\pi[S']$ denotes the marginal distribution over subsequences on $S'$.
\end{msquote}

% ---------------------------------------------------------------- R2.6
\item Could the authors clarify what `register' means in this ribosome-binding
context?

\begin{response}
We have added a clause defining the concept of register:
\end{response}

\begin{msquote}
Doing so, they observed that the SD landscape contains multiple prominent
fitness peaks corresponding to the canonical AGGAG motif appearing in different
registers (i.e., starting at different positions upstream of the start codon).
\end{msquote}

\begin{response}
The following paragraph also carries a footnote stating the indexing convention
used throughout this section:
\end{response}

\begin{msquote}
[FOOTNOTE: Throughout this section we index sequence positions $l$ by their
offset from the start codon, so that $l$ runs from $-13$ to $-5$ rather than from
1 to 9. We also denote by $i$ the register of the AGGAG core motif, so that $i$
runs from $-13$ to $-9$.]
\end{msquote}

% ---------------------------------------------------------------- R2.7
\item Discussion: `it can also be applied to nonlinear and nonparametric models,
such as neural networks or Gaussian processes. This is done by representing the
predicted landscape using an all-order model and applying GaugeFixer to its
parameters.' -- Could the authors unpack this? Does it mean to locally approximate
the deep-learning landscape by an all-order model, a bit like LIME does for linear
regression?

\begin{response}
This procedure does not involve any local approximation, such as that used by
LIME. Rather, one first evaluates the nonlinear or nonparametric model on every
possible sequence. One then expresses the resulting landscape as an all-order
model and uses GaugeFixer to transform its parameters into the desired gauge. We
have revised Results to clarify this point:
\end{response}

\begin{msquote}
Here we used GaugeFixer to characterize the landscape around each fitness peak.
We represented the combinatorially complete landscape of \citet{MartiGomez2026}
as an all-order model having $(4+1)^9 = 1{,}953{,}125$ parameters expressed in
the trivial gauge (in which only the $4^9$ parameters of order 9 are nonzero).
\end{msquote}

\begin{response}
We have revised the Discussion to emphasize this point:
\end{response}

\begin{msquote}
Although GaugeFixer operates only on specific classes of linear models
(all-order and hierarchical models), it can still be used to characterize
nonparametric or nonlinear models (e.g., Gaussian process models or deep neural
networks). To do this, one first evaluates the model of interest on every
possible sequence. The resulting landscape is then represented as an all-order
model, and GaugeFixer is applied to the corresponding parameters. This is what we
did for the Shine-Dalgarno landscape analyzed above (Figure 3). Because this
approach requires explicitly representing the combinatorially complete
landscape, it can be used only when sequences are short enough for that to be
computationally feasible. We note that recent theoretical results, however,
suggest a way of overcoming this limitation for certain Gaussian process models
\citep{Petti2025aa}.
\end{msquote}

% ---------------------------------------------------------------- R2.8
\item Figure 2: could the authors clarify: if a time limit was set for standard
gauge fixing that explains the curves stopping at $10^4$ parameters, and which
one; the characteristics of the model ($\alpha$, $L$, etc.)

\begin{response}
No time limit was imposed on the standard gauge-fixing procedure. Instead, the
benchmark was limited by the memory required to construct the dense projection
matrix. Peak memory reached 1.4 GB for the largest all-order model tested
($9{,}261$ parameters) and 0.6 GB for the largest pairwise model tested
($6{,}121$ parameters). Because this requirement grows as $M^2$, the next
all-order model in the series ($194{,}481$ parameters) would have required
roughly 600 GB, and the next pairwise model ($11{,}361$ parameters) roughly 2 GB.
We therefore restricted the benchmark to tasks requiring less than 2 GB of
memory. We have also stated the range of sequence lengths tested for each model
class. The revised Figure 2 legend reads:
\end{response}

\begin{msquote}
Benchmarking results. Shown are the runtime (A) and peak memory (B) requirements
for GaugeFixer and for standard matrix multiplication, plotted as a function of
the number of model parameters. Tests were carried out for all-order models
($L=2$ to $5$) and pairwise models ($L=3$ to $150$) on protein sequences
($\alpha=20$). Each
data point represents the mean and standard deviation over 10 gauge-fixing
computations performed on a standard laptop computer. No execution-time limit was
imposed, but benchmarking computations for standard matrix multiplication were
restricted to models whose projection matrix required less than 2 GB of memory.
\end{msquote}

% ---------------------------------------------------------------- R2.9
\item Figure 3C: what do the black dots indicate?

\begin{response}
The black dots indicate the fixed reference nucleotides for the core motif in
each register. We have added the following clarification to the figure legend:
\end{response}

\begin{msquote}
Black dots mark the reference nucleotide at each position of the core motif, as
indicated in panel A.
\end{msquote}

% ---------------------------------------------------------------- R2.10
\item Section on: Application to the Shine-Dalgarno fitness landscape. I
understand the authors refer to another work for details, but it would be useful
to report quickly a few details on the parameters shown. How were they inferred?
Was any regularisation used? How many data were used?

\begin{response}
We have expanded the text to state the amount of data used and to describe the
inference procedure. The measured data are now quantified as:
\end{response}

\begin{msquote}
This yielded activity measurements for $257{,}565$ of the $4^9 = 262{,}144$
possible 9-nt sequences.
\end{msquote}

\begin{response}
and the inference procedure is described in the following paragraph:
\end{response}

\begin{msquote}
Recently, \citet{MartiGomez2026} extended this landscape to all possible 9-nt sequences using variance component regression, a Bayesian inference method based on Gaussian processes \citep{Zhou2022}. This method estimates the variance attributable to genetic
interactions of each order, then uses these estimates to define a prior over
landscapes.
\end{msquote}

\begin{response}
Regularization enters through this prior rather than through an explicit penalty
term.
\end{response}

% ---------------------------------------------------------------- R2.11
\item The sentence: `For each peak we define a distribution $\pi$ that fixes the
AGGAG core motif in the corresponding register': could one report the
mathematical formulation of this choice? I think it would be useful to
potentially make these findings relatable to the standard gauges in the field,
i.e., is this choice similar to a zero-sum gauge or others?

\begin{response}
We have revised the main text to provide an explicit mathematical definition of the
distribution associated with each register:
\end{response}

\begin{msquote}
Then for each register $i$, we defined a distribution $\pi$ that fixes the core
motif $u=\textrm{AGGAG}$ at that register while randomizing the other positions
within the 9-nt variable region (Figure 3A). [FOOTNOTE: Formally,
$\pi_{i+j-1}^{c}=x_j^c(u)$ for $1\leq j \leq 5$, $\pi_l^c=0.25$ for $l<i$ or
$l>i+4$.]
\end{msquote}

\begin{response}
The gauges used in this example belong to the family of generalized wild-type
gauges, and we have revised the text to say so explicitly. The interpretation of
the resulting parameters is stated where the gauges are defined, and is then
developed for each parameter type in turn as the corresponding results are
presented (constant parameters, additive parameters, and pairwise interactions,
respectively).
\end{response}

\begin{msquote}
These gauges are examples of the generalized wild-type gauges described
by~\citet{Posfai2025aa}, in which a reference sequence is specified at only a
subset of positions. Under such a gauge, the parameters describe the average
effects of mutations away from the reference sequence, computed over random
backgrounds at the remaining positions. We emphasize that, although the
resulting parameter values differ across gauges, they are alternative
representations of the same underlying model and therefore yield identical
predictions.
\end{msquote}

\end{enumerate}

% =============================================================================
% References cited in this response. Uses the same .bib file as the manuscript.
\bibliographystyle{../kinneylab_bibstyle}
\bibliography{../reference}

\end{document}
